\section{Finite Sets}

\begin{theorem} \label{an:thm:bijection-equivalence-relation}
  Define $\sim$ by
  \[
    \forall X \forall Y, X \sim Y \iff \exists f: X \xrightarrow{\sim} Y
  \]
  Then $\sim$ is an equivalence relation.
\end{theorem}

\begin{proof}
  By \cref{an:def:equivalence-relation},
  \begin{itemize}
    \item \textbf{Reflexivity}: Define $f: X \to X$ by
          \[
            \forall x \in X, f(x) = x
          \]
          Then
          \begin{itemize}
            \item \textbf{Injectivity}: Suppose
                  \begin{itemize}
                    \item $x_{1}, x_{2} \in X$
                    \item $f(x_{1}) = f(x_{2})$
                  \end{itemize}
                  Then
                  \[
                    x_{1} = x_{2}
                  \]
            \item \textbf{Surjectivity}: Suppose $y \in X$. Then
                  \[
                    f(y) = y
                  \]
          \end{itemize}
          Therefore,
          \[
            \begin{array}{c}
              f: X \xrightarrow{\sim} X \\
              \Downarrow                \\
              X \sim X
            \end{array}
          \]
    \item \textbf{Symmetry}: Suppose $X \sim Y$. Then
          \[
            \begin{array}{c}
              \exists f: X \xrightarrow{\sim} Y                           \\
              \by{\Cref{an:thm:inverse-function-is-bijective}} \Downarrow \\
              f^{-1}: Y \xrightarrow{\sim} X                              \\
              \Downarrow                                                  \\
              Y \sim X
            \end{array}
          \]
    \item \textbf{Transitivity}: Suppose
          \begin{itemize}
            \item $X \sim Y$
            \item $Y \sim Z$
          \end{itemize}
          Then
          \begin{itemize}
            \item $\exists f: X \xrightarrow{\sim} Y$
            \item $\exists g: Y \xrightarrow{\sim} Z$
          \end{itemize}
          By \cref{an:thm:properties-of-composition-functions},
          \[
            \begin{array}{c}
              g \circ f: X \xrightarrow{\sim} Z \\
              \Downarrow                        \\
              X \sim Z
            \end{array}
          \]
  \end{itemize}
\end{proof}

\bigskip

\begin{remark}
  Since there is no \textbf{set of all sets} in the underlying set theory, the relation defined in
  \cref{an:thm:bijection-equivalence-relation} is a \textbf{proper-class relation} rather than a set. Its
  restriction to any set $S$ becomes an equivalence relation on $S$.
\end{remark}

\bigskip

\begin{definition} \label{an:def:finite-and-countable-sets}
  Recall the relation $\sim$ defined in \cref{an:thm:bijection-equivalence-relation}.
  A set $X$ is said to be \textbf{finite} if
  \[
    \exists n \in [0, \infty)_{\N}, [1, n]_{\N} \sim X
  \]
  Such an $n$ is called a \textbf{cardinality} of $X$, and is denoted by
  \[
    \abs*{X} = n
  \]
  The notation $\abs*{X} < \aleph_{0}$ denotes that $X$ is finite.
\end{definition}

\bigskip

\begin{definition} \label{an:def:infinite-countable-uncountable-sets}
  A set $X$ is said to be
  \begin{itemize}
    \item \textbf{infinite} if it is not finite
    \item \textbf{countable} if $\N \sim X$
    \item \textbf{at most countable} if it is either finite or countable
    \item \textbf{uncountable} if it is neither finite nor countable
  \end{itemize}
  The notation
  \begin{itemize}
    \item $\abs*{X} \geq \aleph_{0}$ denotes that $X$ is infinite
    \item $\abs*{X} = \aleph_{0}$ denotes that $X$ is countable
    \item $\abs*{X} \leq \aleph_{0}$ denotes that $X$ is at most countable
    \item $\abs*{X} > \aleph_{0}$ denotes that $X$ is uncountable
  \end{itemize}
  The order symbols above are not \textbf{total orders}; they are just convenient notations.
\end{definition}

\bigskip

\begin{lemma} \label{an:lem:subset-of-Nn-is-finite}
  \[
    \forall n \in [0, \infty)_{\N}, \forall A \subseteq [1, n]_{\N}, \abs*{A} < \aleph_{0}
  \]
\end{lemma}

\begin{proof}
  Define
  \[
    \forall n \in [0, \infty)_{\N}, P(n): \forall A \subseteq [1, n]_{\N}, \abs*{A} < \aleph_{0}
  \]
  Then
  \begin{itemize}
    \item \textbf{Base case:} Define $f: [1, 0]_{\N} \xrightarrow{\sim} [1, 0]_{\N}$ by
          \[
            \forall k \in [1, 0]_{\N}, f(k) = k
          \]
          Then
          \[
            \abs*{[1, 0]_{\N}} = 0 < \aleph_{0}
          \]
          Therefore,
          \[
            \begin{array}{c}
              \forall A \subseteq [1, 0]_{\N}, A = [1, 0]_{\N}           \\
              \Downarrow                                                 \\
              \forall A \subseteq [1, 0]_{\N}, \abs*{A} = 0 < \aleph_{0} \\
              \Downarrow                                                 \\
              P(0) \text{ is true}
            \end{array}
          \]
    \item \textbf{Inductive step:} Suppose
          \begin{itemize}
            \item $P(n)$ be true
            \item $A \subseteq [1, n + 1]_{\N}$
          \end{itemize}
          Then
          \begin{itemize}
            \item If $n + 1 \notin A$, then
                  \[
                    \begin{array}{c}
                      A \subseteq [1, n]_{\N}                   \\
                      \by{\ensuremath{P(n)} is true} \Downarrow \\
                      \abs*{A} < \aleph_{0}                     \\
                    \end{array}
                  \]
            \item Otherwise, define
                  \[
                    B \coloneqq A \setminus \set*{n + 1} \subseteq [1, n]_{\N}
                  \]
                  Then
                  \[
                    \begin{array}{c}
                      B \subseteq [1, n]_{\N}                   \\
                      \by{\ensuremath{P(n)} is true} \Downarrow \\
                      \abs*{B} < \aleph_{0}                     \\
                      \Downarrow                                \\
                      \exists m \in [0, \infty)_{\N}, \exists f: [1, m]_{\N} \xrightarrow{\sim} B
                    \end{array}
                  \]
                  Define $g: [1, m + 1]_{\N} \to A$ by
                  \[
                    \forall k \in [1, m + 1]_{\N}, g(k) = \begin{dcases}
                      n + 1, & \text{if } k = m + 1 \\
                      f(k),  & \text{otherwise}
                    \end{dcases}
                  \]
                  Then
                  \begin{itemize}
                    \item \textbf{Injectivity}: Suppose
                          \begin{itemize}
                            \item $k_{1}, k_{2} \in [1, m + 1]_{\N}$
                            \item $g(k_{1}) = g(k_{2})$
                          \end{itemize}
                          Then
                          \begin{itemize}
                            \item If $k_{1} = m + 1$, then
                                  \[
                                    \begin{array}{c}
                                      g(k_{1}) = n + 1                            \\
                                      \Downarrow                                  \\
                                      g(k_{2}) = n + 1                            \\
                                      \by{\ensuremath{n + 1 \notin B}} \Downarrow \\
                                      k_{2} = m + 1
                                    \end{array}
                                  \]
                            \item Otherwise,
                                  \[
                                    \begin{array}{c}
                                      k_{1} \in [1, m]_{\N}                                            \\
                                      \Downarrow                                                       \\
                                      g(k_{1}) = f(k_{1})                                              \\
                                      \Downarrow                                                       \\
                                      g(k_{2}) = f(k_{1})                                              \\
                                      \by{\ensuremath{f(k_{1}) \neq n + 1}} \Downarrow                 \\
                                      g(k_{2}) = f(k_{2}) = f(k_{1})                                   \\
                                      \by{\ensuremath{f: [1, m]_{\N} \xrightarrow{\sim} B}} \Downarrow \\
                                      k_{1} = k_{2}
                                    \end{array}
                                  \]
                          \end{itemize}
                    \item \textbf{Surjectivity}: Suppose $a \in A$.
                          \begin{itemize}
                            \item If $a = n + 1$, then
                                  \[
                                    g(m + 1) = a
                                  \]
                            \item Otherwise,
                                  \[
                                    \begin{array}{c}
                                      a \in B                                                          \\
                                      \by{\ensuremath{f: [1, m]_{\N} \xrightarrow{\sim} B}} \Downarrow \\
                                      \exists k \in [1, m]_{\N}, f(k) = a                              \\
                                      \Downarrow                                                       \\
                                      g(k) = f(k) = a
                                    \end{array}
                                  \]
                          \end{itemize}
                  \end{itemize}
                  Therefore,
                  \[
                    \begin{array}{c}
                      g: [1, m + 1]_{\N} \xrightarrow{\sim} A \\
                      \Downarrow                              \\
                      \abs*{A} = m + 1 < \aleph_{0}           \\
                    \end{array}
                  \]
          \end{itemize}
          Therefore, $P(n + 1)$ is true.
  \end{itemize}
  By \cref{an:def:mathematical-induction},
  \[
    \forall n \in [0, \infty)_{\N}, P(n) \text{ is true}
  \]
\end{proof}

\bigskip

\begin{theorem} \label{an:thm:subset-of-finite-set-is-finite}
  \[
    \forall X \forall Y, \abs*{X} < \aleph_{0} \land Y \subseteq X \implies \abs*{Y} < \aleph_{0}
  \]
\end{theorem}

\begin{proof}
  Suppose
  \begin{itemize}
    \item $\abs*{X} < \aleph_{0}$
    \item $Y \subseteq X$
  \end{itemize}
  Then
  \[
    \exists n \in [0, \infty)_{\N}, \exists f: [1, n]_{\N} \xrightarrow{\sim} X
  \]
  Define
  \[
    Z \coloneqq f^{-1}(Y) \subseteq [1, n]_{\N}
  \]
  By \cref{an:lem:subset-of-Nn-is-finite},
  \[
    \begin{array}{c}
      \abs*{Z} < \aleph_{0} \\
      \Downarrow            \\
      \exists m \in [0, \infty)_{\N}, [1, m]_{\N} \sim Z
    \end{array}
  \]
  By \cref{an:thm:properties-of-restriction-functions},
  \[
    \begin{array}{c}
      f|_{Z}^{Y}: Z \xrightarrow{\sim} Y                           \\
      \by{\Cref{an:thm:bijection-equivalence-relation}} \Downarrow \\
      \abs*{Y} = m < \aleph_{0}
    \end{array}
  \]
\end{proof}
